\documentclass[11pt,a4paper]{article}

% ---------------------------------------------------------------------------
% Korean encoding tax: a prospectively frozen code-length measurement
% Draft generated 2026-09-10 from the frozen measurement spec (hash 850fa9be)
% and the held-out analysis in pipeline/results/. Every number in the tables
% is copied from results/*/analysis_summary.json; the tokenizer-generation
% counterfactual in Sec. 5.4 is post hoc and labelled as such.
% ---------------------------------------------------------------------------

\usepackage[utf8]{inputenc}
\usepackage[T1]{fontenc}
\usepackage{lmodern}
\usepackage[margin=2.5cm]{geometry}
\usepackage{amsmath,amssymb}
\usepackage{booktabs}
\usepackage{array}
\usepackage{tabularx}
\usepackage{float}
\usepackage[section]{placeins}
\usepackage{microtype}
\usepackage{hyperref}
\usepackage{xcolor}
\usepackage[numbers,sort&compress]{natbib}
\usepackage{caption}
\usepackage{enumitem}
\usepackage{needspace}

\hypersetup{
  colorlinks=true,
  linkcolor=blue!60!black,
  citecolor=blue!60!black,
  urlcolor=blue!60!black,
  pdftitle={Is Korean's Higher LLM Token Cost Information or Encoding?},
  pdfauthor={Jaehyung Park},
  pdfsubject={A prospectively frozen Korean--English tokenization study},
  pdfkeywords={Korean, English, tokenization, code length, language models}
}
\captionsetup{font=small,labelfont=bf}
\renewcommand{\bibfont}{\small}

\newcommand{\RT}{R_{T}}
\newcommand{\RB}{R_{B}}
\newcommand{\Reta}{R_{\eta}}
\newcommand{\CK}{C_{K}}
\newcommand{\ko}{\mathrm{KR}}
\newcommand{\en}{\mathrm{EN}}

\title{Is Korean's Higher LLM Token Cost Information or Encoding?\\
\large A prospectively frozen code-length measurement on an aligned Korean--English corpus}

\author{Jaehyung Park\\
Department of Electrical and Computer Engineering, Seoul National University\\
\texttt{jh0420park@gmail.com}}
% TODO: confirm the romanisation of 박재형 before submission.

\date{Revised draft of 11 September 2026}

\begin{document}
\maketitle

\begin{abstract}
Aligned Korean text often consumes more billed tokens than English text with
parallel content. Two contributions may coexist: the texts can differ in
predictive code length, and a billing tokenizer can represent them with
different efficiency. Token counts alone cannot distinguish these
contributions. We operationalise code length with a fixed language model,
$B_M(s)=-\sum_j\log_2 Q_M(z_j\mid z_{<j})$, and compare its corpus sum with
content-only token counts. On 7{,}046 aligned sentence pairs from 278 Global
Voices articles, analysed after the held-out plan was prospectively frozen,
the Korean side used 1.389 times as many \texttt{o200k\_base} tokens as the
English side (95\% cluster-bootstrap CI 1.371--1.406) and had 1.312 times the
Qwen3-1.7B-Base-relative code length (1.296--1.328). Their ratio,
$\Reta=\RB/\RT=0.945$ (0.936--0.953), corresponds to a 5.5\% lower
model-relative code length per billed Korean token. The primary rule required
the one-sided 95\% upper bound of $\Reta$ to be below 0.95; the observed bound
was 0.952, so both the primary rule and the overall all-pass rule failed. The
token-ratio criterion passed, and a separate tokenizer-sensitivity control
passed: Polyglot-Ko assigned the identical Korean strings 13\% fewer token
units ($\CK=0.867$). Exploratory bit-ratio equivalence failed, contradicting a
simple encoding-only pattern in which total code lengths are equal. The study
therefore demonstrates token-count dependence on tokenizer choice but neither
establishes the prespecified 5\% efficiency shortfall nor identifies a causal
division between linguistic, translation, model, and tokenizer effects.
\end{abstract}

\section{Introduction}

Users of LLM application programming interfaces (APIs) are commonly charged
per token. Parallel-language studies have shown that the same content can
require substantially different token counts, producing unequal monetary and
context-window costs \citep{ahia2023,petrov2023}. Training-corpus composition
is one plausible contributor, but it is not the only one. Korean and English
versions are translations rather than identical strings; a translation can
expand or contract, and a fixed model can predict the two languages with
different accuracy. Token count alone cannot identify these mechanisms.

An information-theoretic comparison puts the two measurements on explicit
scales: tokenizer symbols and code bits under a stated predictive model. It
does not, by itself, recover intrinsic semantic information or a causal
decomposition. If token and model-relative bit ratios move together, the
billing-token surplus is accompanied by a similar predictive-code-length
surplus. If they separate, tokenizer efficiency is one possible contributor,
subject to model and translation confounding.

\paragraph{Related work.} Cross-lingual token-count and billing disparities
are established empirically \citep{ahia2023,petrov2023}. Information-theoretic
work has evaluated tokenizers through channel use, entropy, and compression
\citep{zouhar2023,erdogan2026}. A recent multilingual token-cost ledger
separates removable coding terms from proposed intrinsic terms on FLORES-200
\citep{mandarapu2026}. The narrower contribution here is a Korean--English
case study that combines production-token counts with Qwen-relative
sequential code length under a prospectively frozen held-out analysis; it
does not claim the first general decomposition or an estimator of intrinsic
language entropy.

This paper reports an operational comparison with decision rules fixed in
advance. The contributions are: (i) a pooled efficiency ratio $\Reta$ that
compares Qwen-relative code length with billed-token use; (ii) a
pipeline with unit tests that computes it reproducibly with pinned model,
tokenizer, and corpus revisions; (iii) a pilot/held-out split with
thresholds frozen by content hash before the held-out data were scored; and
(iv) the held-out result, which narrowly fails the prespecified primary rule,
together with clearly labelled exploratory analyses.

\section{Background: tokenizers as codes, models as codebooks}
\label{sec:background}

Shannon's source-coding theorem distinguishes the information in a message
from the number of symbols used to write it down \citep{shannon1948}. The
ideal fractional code length for an event of probability $p$ is
$-\log_2 p$ bits; literal prefix-code words have integer lengths, while block
and arithmetic coding approach the ideal average.
Shannon's later experiment on printed English estimated the entropy of a
language from a human predictor's guesses \citep{shannon1951}; a language
model plays the same predictive role today, with the difference that its
predictions are explicit probabilities.

Let $M$ be a fixed autoregressive language model with its own tokenizer
$\tau_M$. For a message $s$, write $z=\tau_M(s)=(z_1,\dots,z_n)$ and let
$c$ be the fixed document-boundary token. The \emph{model-relative code
length} is
\begin{equation}
  B_M(s) \;=\; -\sum_{j=1}^{n} \log_2 Q_M\!\left(z_j \mid c,z_{<j}\right),
  \label{eq:bits}
\end{equation}
the ideal sequential code length assigned by $M$; an arithmetic coder driven
by these probabilities approaches this length up to a small finite-coding
overhead \citep{mackay2003}. No end-of-sequence event is scored, so this is a
conditional code length whose decoder is assumed to know the message boundary
or content-token count $n$. Two properties matter. First, $B_M(s)$ does not
depend on the \emph{billing} tokenizer: it is computed through $\tau_M$ only.
Second, it is model-relative surprisal, not intrinsic semantic information.
Writing $P$ for the true conditional
distribution and $Q$ for the model's, the expected code length decomposes as
$H(P,Q)=H(P)+D_{\mathrm{KL}}(P\,\|\,Q)$, so $B_M$ includes the model's own
prediction error. A model that knows Korean less well than English will
overprice Korean text; the direction of this bias is discussed in
Section~\ref{sec:discussion}.

For a billing tokenizer $b$ with $T_b(s)=|\tau_b(s)|$, its operational
efficiency on $s$ is the model-relative code bits per billed token,
$\eta_{M,b}(s)=B_M(s)/T_b(s)$. The corpus-level estimands are ratios of
sums over aligned pairs $i$:
\begin{equation}
  \RT=\frac{\sum_i T_{i,\ko}}{\sum_i T_{i,\en}},\qquad
  \RB=\frac{\sum_i B_{i,\ko}}{\sum_i B_{i,\en}},\qquad
  \Reta=\frac{\sum_i B_{i,\ko}/\sum_i T_{i,\ko}}{\sum_i B_{i,\en}/\sum_i T_{i,\en}}
       =\frac{\RB}{\RT}.
  \label{eq:estimands}
\end{equation}
As simplified diagnostic patterns, an information-only account predicts
$\RB\approx\RT$ and $\Reta\approx 1$, whereas an encoding-only account
predicts $\RT>1$, $\RB\approx 1$, and $\Reta<1$. Real data can contain both
contributions, and $B_M$ also contains model mismatch. The quantity
$1-\Reta$ is the relative shortfall in model-relative code bits per Korean
token compared with an English-efficiency counterfactual. It is not a literal count
of empty tokens or money that purchased no information. Ratios of sums define
a token-weighted corpus estimand and preserve $\Reta=\RB/\RT$. Means of
per-message ratios answer a different, equally-pair-weighted question and can
be unstable when denominators are small.

An independent control uses a tokenizer trained on Korean, $k$. The Korean
strings and their Qwen-relative code lengths are unchanged; only the counting
tokenizer changes.
The absolute Korean-token reduction
\begin{equation}
  \CK=\frac{\sum_i T^{(k)}_{i,\ko}}{\sum_i T^{(b)}_{i,\ko}}
\end{equation}
measures the token-count change for identical Korean strings under the
comparison code. $\CK<1$ shows tokenizer dependence, but it is not a directly
deployable billing counterfactual because a model's tokenizer cannot normally
be replaced without retraining or adaptation. Attributing the difference to
training-language balance alone would also require holding the algorithm,
vocabulary size, normalisation, and pre-tokenisation fixed, which these two
off-the-shelf tokenizers do not.

For completeness, the exploratory relative gap-contraction statistic is
$G=\RT^{(k)}/\RT^{(b)}$. It can decrease because the Korean/English ratio
narrows by either fewer Korean tokens or more English tokens, so $\CK$ is the
direct absolute Korean-token control.

\section{Measurement design}
\label{sec:design}

\subsection{Corpus and split}

We use the OPUS packaging of Global Voices v2018q4 in English and Korean
\citep{tiedemann2012opus,prokopidis2016globalvoices}, a citizen-media news collection whose articles are
translated by volunteers. The archive contains 365 article-level alignment
groups, of which 347 contain aligned text, yielding 9{,}017 aligned sentence
pairs. Article identity is preserved as a cluster identifier; the corpus has
one genre (news), and the translation direction is not recorded.

The release has no train/test division, so a deterministic article-level
split was created before any scoring. Articles connected by an identical
bilingual pair were merged into indivisible components; components were
ranked by a seeded SHA-256 rule (seed 20260910) and the prefix closest to 20\%
of articles was assigned to the pilot. The pilot contains 69 articles
(1{,}732 pairs) and the held-out set 278 articles (7{,}285 pairs); no article
or exact bilingual pair crosses the split. The files were created, hashed,
and mechanically validated during corpus preparation, but no held-out outcome
measurements or summaries were examined before the analysis specification was
frozen.

Only Unicode NFC normalisation and surrounding-whitespace removal are
applied. Validation flags record rows whose Korean side contains no Hangul,
rows whose Korean and English text are identical (untranslated photo credits
and licence strings), disallowed control characters, and length statistics.

\subsection{Estimator model and scoring policy}

The estimator is Qwen3-1.7B-Base \citep{qwen3,qwen3model} at revision
\texttt{ea980cb0}, run in bfloat16 with the log-softmax computed in
float32. Every message is scored in evaluation and inference mode from raw
next-token logits. The model's document-boundary token
(\texttt{<|endoftext|>}, id 151643) is prepended and every content token is
scored conditional on it, so that the first content token receives a code
length; no end-of-sequence event is charged. Special tokens are never added
implicitly (\texttt{add\_special\_tokens=False}), padding is never scored,
and the pipeline asserts on every call that the number of scored targets
equals the number of content tokens. The observed maxima were 185 Korean and
126 English estimator tokens, far below the 32{,}768-token context, so no
chunking occurred.

\subsection{Tokenizers}

The baseline billing tokenizer is OpenAI's \texttt{o200k\_base} encoding
(tiktoken 0.14.0; 199{,}998 mergeable tokens), used by the GPT-4o
generation \citep{tiktoken}. The Korean-trained comparison tokenizer is that
of EleutherAI's Polyglot-Ko-1.3B \citep{polyglotko,polyglotkoartifact} at revision
\texttt{557e162c}, a 30{,}000-entry byte-level BPE trained on a Korean-only
corpus. Token counts are content-only, with no chat template or wrapper
tokens. Counting tokenizers return integers only; their token ids are never
fed to the estimator, and the estimator is scored only through its own
tokenizer. The estimator snapshot and baseline tokenizer asset are
hash-verified before scoring. The comparison tokenizer is revision-pinned;
its resolved files, like those of all three components, are hashed in the run
provenance.

\subsection{Prospectively frozen analysis}
\label{sec:prereg}

Ratios described as ``well above one'' or ``close to one'' are meaningless
until thresholds are fixed, and thresholds chosen after seeing the data
prove nothing. All thresholds, endpoints, and exclusion rules were therefore
written into a specification file whose content hash (SHA-256
\texttt{850fa9be}\ldots) was recorded when it was frozen on 10 September
2026, after the pilot had been analysed and before the held-out set was
scored. A pre-results Git commit (\texttt{2a6d202}) provides an additional
time-ordered record, but the plan was not deposited in an external
pre-registration registry. Confirmatory analysis validates the frozen content
hash and refuses pilot-tagged inputs; these internal checks reduce accidental
deviation but are not an external cryptographic registration. The frozen
decisions are listed in Table~\ref{tab:spec}.

\begin{table}[htbp]
\centering\small
\caption{Frozen decisions. Bounds are cluster-bootstrap percentiles.}
\label{tab:spec}
\begin{tabularx}{\textwidth}{@{}l>{\raggedright\arraybackslash}p{5.2cm}>{\raggedright\arraybackslash}X@{}}
\toprule
Decision & Value & Rule on the held-out bootstrap \\
\midrule
Primary endpoint & $\Reta$ & one-sided 95\% upper bound $<c_\eta=0.95$ \\
Token ratio & $\RT$ & one-sided 95\% lower bound $>c_T=1.20$ \\
Tokenizer control & $\CK$ & one-sided 95\% upper bound $<c_C=0.90$ \\
Overall rule & all three confirmatory comparisons must pass & global conjunction: all three pass \\
Bit-ratio equivalence (exploratory) & $\RB$ & 90\% CI inside $[0.8333,\,1.20]$ \\
Exclusions & identical KR/EN text; no Hangul on KR side & applied identically to pilot and held-out \\
Bootstrap & 10{,}000 replicates, seed 20260910 & whole articles resampled, both languages together \\
\bottomrule
\end{tabularx}
\end{table}

The margins operationalise effects judged practically material for this
study: at least 20\% extra billed tokens ($c_T=1.20$), at least a 5\%
efficiency shortfall ($c_\eta=0.95$), and at least a 10\% Korean-token
reduction under the comparison tokenizer ($c_C=0.90$). They are
researcher-chosen margins, not externally standardised cutoffs. The
exploratory bit-ratio equivalence interval allows a reciprocal 20\% departure
from equality. The pilot informed feasibility and precision. Its point estimates
($\Reta=0.940$, $\RT=1.377$, $\CK=0.870$ after exclusions) located the
effect, and a precision simulation, which resampled the pilot's per-article
sums to the held-out article count, described the distribution of the
bootstrap bounds that the rules would use. For $c_\eta=0.95$ that
simulation gave a pass probability of 0.78 conditional on the pilot values
carrying over unchanged; the threshold lay 1.25 pilot standard errors above
the pilot point estimate. This is recorded here because it bears directly on
the interpretation of the confirmatory outcome.

The frozen confirmatory claim was the conjunction $\Reta<0.95$,
$\RT>1.20$, and $\CK<0.90$ under their stated one-sided bounds. Because
$\RB$ equivalence was exploratory, this design does not confirmatorily
adjudicate the full simple encoding-only prediction $\RB\approx1$.

\subsection{Statistical analysis}

All estimands are ratios of corpus sums, Eq.~\eqref{eq:estimands}.
Uncertainty is quantified by a paired, stratified cluster bootstrap:
articles are treated as exchangeable draws and resampled with replacement
within genre strata, every pair of a
resampled article enters with both its Korean and English measurements, and
each pooled ratio is recomputed in every replicate. Because the estimands are
functions of column sums only, per-article sums are precomputed and the
result is computationally identical to row-level resampling of whole
articles. The resulting intervals target a superpopulation of Global
Voices-style Korean--English volunteer-translated citizen-media/news
articles, not all Korean and English use. The identity $\Reta=\RB/\RT$ is asserted in every
replicate. One-sided decision bounds are the 5th or 95th percentiles;
equivalence for $\RB$ uses the two one-sided tests procedure, requiring the
90\% percentile interval to lie inside the pre-specified interval
\citep{schuirmann1987}. The all-pass rule is one global conjunction; every
component must pass, so no multiplicity correction is invoked for that global
claim.

Exploratory analyses, which never feed the confirmatory rule, are: pooled
estimands within terciles of English character count; the distribution of
per-message ratios; and paired cluster-level sign-flip permutation tests in
which the Korean and English measurements of all pairs in a randomly chosen
set of articles are exchanged, with the log pooled ratio as statistic and
10{,}000 permutations. These tests address the sharp null that language
labels (or, for $\CK$, tokenizer labels) are exchangeable within article,
which is stronger than the statement that a pooled ratio equals one. The
three language ratios are functionally dependent and are reported as one
finding seen from three angles.

\subsection{Implementation}

The pipeline is a Python package (\texttt{klen}) with 66 unit tests,
including hand-verifiable code-length cases (a uniform fake model yields
exactly $\log_2 10$ bits per token; a biased fake model exactly 1 bit),
whole-cluster resampling verified against the enumerable set of possible
replicates, stratum preservation, freeze and tamper detection, and an
independent re-implementation cross-check against a tiny random model.
Provenance for every run records package versions (PyTorch 2.14.0,
Transformers 5.17.0, Python 3.13.1), tokenizer file hashes, the dataset
hash, the platform (\texttt{arm64}) and device backend (\texttt{mps}), the
specification, and the validation report.

\Needspace{10\baselineskip}
\section{Results}
\label{sec:results}

\subsection{Sample}

After the frozen exclusion rules, the held-out set comprises 7{,}046 pairs
in 278 articles (239 rows dropped: 104 with identical Korean and English
text and a further 135 with no Hangul on the Korean side). The pilot after
the same rules comprises 1{,}672 pairs in 69 articles. No row changed under
NFC normalisation; no pair id was duplicated.

\subsection{Pooled estimands and confirmatory decisions}

\begin{table}[htbp]
\centering\small
\caption{Held-out pooled estimands (278 articles, 7{,}046 pairs) with 95\%
cluster-bootstrap percentile intervals and prespecified decisions.
$\RT$ counts \texttt{o200k\_base} tokens. $\CK$ is the
Polyglot-Ko/\texttt{o200k\_base} ratio for the same Korean strings.
Bits are Qwen3-1.7B-Base-relative code lengths.}
\label{tab:main}
\begin{tabularx}{\textwidth}{@{}>{\raggedright\arraybackslash}p{4.6cm}cc>{\raggedright\arraybackslash}Xl@{}}
\toprule
Quantity & Point & 95\% CI & Decision bound & Result \\
\midrule
$\RT$ token ratio KR/EN & 1.3888 & [1.3714, 1.4058] & lower 1.3741 $>1.20$ & pass \\
$\RB$ bit ratio KR/EN & 1.3121 & [1.2959, 1.3278] & 90\% CI [1.2984, 1.3252], outside $[0.8333, 1.20]$ & no equivalence$^\ast$ \\
$\Reta=\RB/\RT$ (primary) & 0.9448 & [0.9363, 0.9532] & upper 0.9519 $\not<0.95$ & \textbf{fail} \\
$1-\Reta$ shortfall & 0.0552 & [0.0468, 0.0637] & --- & --- \\
$\CK$ Korean tokens, Polyglot/o200k & 0.8665 & [0.8601, 0.8729] & upper 0.8719 $<0.90$ & pass \\
$\RT$ under Polyglot-Ko & 0.5095 & [0.5025, 0.5163] & --- & exploratory \\
$G$ gap contraction & 0.3668 & [0.3629, 0.3707] & --- & exploratory \\
\midrule
Overall (all three criteria) & & & & \textbf{fail} \\
\bottomrule
\end{tabularx}
\par\smallskip
\raggedright\footnotesize $^\ast$Exploratory only; not part of the overall
confirmatory rule.
\end{table}

Table~\ref{tab:main} gives the held-out estimates. The Korean side costs
38.9\% more \texttt{o200k\_base} tokens and has 31.2\% greater
Qwen-relative code length than the English side. The efficiency ratio is
$\Reta=0.945$, corresponding to 5.5\% lower Qwen-relative code bits per
billed Korean token, with a 95\% interval of 4.7\% to 6.4\%. The
token-ratio and tokenizer-sensitivity component rules pass their prespecified
bounds. The primary comparison does not:
the one-sided 95\% upper bound of $\Reta$ is 0.9519, above the frozen
threshold of 0.95, and under the overall rule the confirmatory claim is
therefore not supported. The bit ratio is outside the exploratory
equivalence interval; the simple encoding-only pattern $\RB\approx1$ is not
observed.

The pilot and held-out estimates agree closely (Table~\ref{tab:pilot}). The
held-out $\Reta$ is 0.0048 above the pilot value, about 1.1 held-out
bootstrap standard deviations; with the threshold placed 1.25
pilot standard errors above the pilot estimate, that shift was enough.

\begin{table}[htbp]
\centering\small
\caption{Pilot versus held-out point estimates after identical exclusions.}
\label{tab:pilot}
\begin{tabular}{@{}lcc@{}}
\toprule
Quantity & Pilot (69 articles, 1{,}672 pairs) & Held-out (278 articles, 7{,}046 pairs) \\
\midrule
$\RT$   & 1.3769 & 1.3888 \\
$\RB$   & 1.2943 & 1.3121 \\
$\Reta$ & 0.9400 & 0.9448 \\
$\CK$   & 0.8697 & 0.8665 \\
\bottomrule
\end{tabular}
\end{table}

\subsection{Exploratory analyses}

\paragraph{Message length.} Table~\ref{tab:length} splits the held-out pairs
into terciles of English character count. The token surplus, the bit
surplus, and the efficiency shortfall all shrink with length: $1-\Reta$ is
9.3\% in the shortest tercile and 4.4\% in the longest; their marginal
intervals do not overlap, but no direct between-tercile contrast was tested.
Because pooled ratios weight by tokens, the corpus-level estimate is
dominated by long messages. Each message includes one boundary-conditioned
first-token term, which occupies a larger share of a short message, but no
alternative-BOS sensitivity analysis was run; the cause of the pattern is
therefore unresolved.

\begin{table}[htbp]
\centering\small
\caption{Held-out pooled estimands within terciles of English character
count (exploratory). Articles may span terciles; bootstrap within tercile.}
\label{tab:length}
\begin{tabular}{@{}lcccccc@{}}
\toprule
Tercile & Chars (EN) & Pairs & $\RT$ & $\RB$ & $\Reta$ [95\% CI] & $\CK$ \\
\midrule
Q1 (short) & $[3,74]$    & 2{,}349 & 1.564 & 1.419 & 0.907 [0.894, 0.921] & 0.882 \\
Q2         & $(74,137]$  & 2{,}348 & 1.415 & 1.318 & 0.931 [0.919, 0.944] & 0.868 \\
Q3 (long)  & $(137,587]$ & 2{,}349 & 1.329 & 1.271 & 0.956 [0.946, 0.966] & 0.861 \\
\bottomrule
\end{tabular}
\end{table}

\paragraph{Per-message distribution.} Per-message ratios are descriptive
only. The median per-pair token ratio is 1.429 (5th--95th percentiles
0.762--2.364), the median bit ratio 1.357 (0.766--2.091), and the median
efficiency ratio 0.944 (0.666--1.369). Their arithmetic means (1.489, 1.396,
0.973) happen to exceed the pooled estimates here. These means are
equal-pair-weighted descriptive summaries, not biased estimates of the
token-weighted pooled quantities.

\paragraph{Permutation tests.} Exchanging the Korean and English
measurements of whole articles 10{,}000 times gives observed log pooled
ratios of $+0.328$ ($\RT$), $+0.272$ ($\RB$), $-0.057$ ($\Reta$) and
$-0.143$ ($\CK$), each with two-sided $p=1/10{,}001$, the smallest
attainable value. Under the sharp label-exchangeability null, every observed
statistic is in the extreme permutation tail. These exploratory tests are not
evidence for the prespecified effect sizes and are not the basis of the
confirmatory decisions.

\subsection{Post hoc: the baseline tokenizer generation}
\label{sec:generation}

To understand why the effect is small, we recomputed $\RT$ and $\Reta$ on
the held-out set holding the Qwen3 bits fixed and substituting the two
previous OpenAI encodings for the billing tokenizer (Table~\ref{tab:gen}).
This was done after the confirmatory analysis and is not part of it.

\begin{table}[htbp]
\centering\small
\caption{Held-out $\RT$ and $\Reta$ under three OpenAI encodings with the
same Qwen3-1.7B-Base bits (post hoc, exploratory).}
\label{tab:gen}
\begin{tabular}{@{}llcc@{}}
\toprule
Encoding & Associated model family & $\RT$ & $\Reta$ \\
\midrule
\texttt{r50k\_base}   & GPT-3            & 4.591 & 0.286 \\
\texttt{cl100k\_base} & GPT-3.5 / GPT-4  & 2.238 & 0.586 \\
\texttt{o200k\_base}  & GPT-4o (frozen)  & 1.389 & 0.945 \\
\bottomrule
\end{tabular}
\end{table}

The model-relative efficiency shortfall is 71\% with
\texttt{r50k\_base}, 41\% with \texttt{cl100k\_base}, and 5.5\% with
\texttt{o200k\_base}. For these same strings, the sequence is consistent
with more efficient Korean token allocation in newer OpenAI encodings. This
post hoc comparison neither identifies vocabulary size or multilingual
training as the cause nor reconstructs historical model or API costs.

\section{Discussion}
\label{sec:discussion}

\subsection{Why the primary comparison failed}

Two distinct things happened. Statistically, the threshold was set 1.25
pilot standard errors above the pilot estimate, the
pre-freeze simulation gave the rule only a 78\% chance of passing even if
the pilot value were exactly right, and the held-out estimate moved up by an
unremarkable 1.1 bootstrap standard deviations. None of this justifies
revising the threshold after observing the result; that constraint is the
purpose of freezing the rule.

Substantively, $\Reta$ is close to one because $\RB$ is large, not because
$\RT$ is small. Under the current tokenizer the Korean token surplus is
39\%, and the model spends 31\% more bits on the Korean side. Three
contributors, which the present design cannot separate, push $\RB$ upward:

\begin{enumerate}[nosep]
\item \emph{Estimator representation.} Bits are summed over the estimator's
own tokens, and Qwen3's tokenizer assigns the Korean side 1.554 times as many
tokens as the English side. Tokenization defines the model's factorisation,
but finer segmentation alone does not mathematically imply a larger total
chain-rule code length. This diagnostic therefore does not identify how much
of $\RB$ comes from the representation.
\item \emph{Model fluency.} The $D_{\mathrm{KL}}$ term is larger for the
language the model predicts less well. Qwen3 reports multilingual training,
but the public model report does not disclose the language proportions
needed to calibrate a Korean--English KL-gap difference. The
held-out code lengths are 4.87 bits per Qwen token for Korean versus 5.77
for English, and 1.42 versus 1.23 bits per UTF-8 byte; both are
unit-dependent diagnostics rather than a calibration of fluency.
\item \emph{Translation content.} Global Voices Korean articles are
volunteer translations that may add explanation; parallel meaning does not
guarantee equal information.
\end{enumerate}

The model-mismatch and translation terms can alter $\RB$, but their
magnitudes and net direction are not measured here. Consequently, the 5.5\%
shortfall is not a lower bound on a causal encoding effect.

\subsection{What the evidence supports}

The prespecified claim of at least a 5\% Korean shortfall in model-relative
code bits per billed token was not established at the 95\% level. Within the
stated corpus and estimators, the evidence supports a 38.9\% Korean
\texttt{o200k\_base} token surplus and passage of its 20\% component
criterion. It also supports tokenizer sensitivity: Polyglot-Ko assigned
13.3\% fewer Korean token units and passed the 10\% component criterion.
However, the exploratory $\RB$ equivalence prediction failed, the primary
$\Reta$ rule failed, and therefore the compound confirmatory hypothesis
failed. The historical-tokenizer pattern is descriptive rather than causal.

\subsection{Limitations}

The estimator is a single small model; $B_M$ is Qwen3-1.7B-relative
surprisal, not semantic information, and its cross-language KL gap is
unquantified. Omitting EOS makes it a code length conditional on external
message framing. Content-only token counts also omit chat wrappers and other
API-specific overhead. A second estimator and the multi-genre FLORES+ corpus
\citep{floresplus}, for which a specification file is prepared, remain
exploratory replications.

The corpus contains one genre with unknown translation direction and predates
the model's release, so training-data contamination cannot be ruled out. The
bootstrap generalises only to Global Voices-style translated articles;
repeated content linking a small number of articles may also make
article-only intervals slightly optimistic. The three language ratios are
functionally dependent. $\CK$ shows token-count sensitivity but neither its
cause nor a tokenizer that can be swapped into an existing billed model.
Historical encodings and $\RB$ equivalence are exploratory. Most importantly,
the present held-out outcomes have now been opened: any rerun or revised
threshold on this corpus is exploratory. A new confirmatory analysis would
require genuinely new, untouched data frozen under a new plan.

\section{Conclusion}

The overall confirmatory hypothesis failed because its primary efficiency
criterion failed. On 7{,}046 held-out aligned pairs, Korean used 38.9\% more
\texttt{o200k\_base} tokens and had 31.2\% greater Qwen3-relative code length,
giving an estimated 5.5\% efficiency shortfall (95\% CI 4.7--6.4\%). The
prespecified rule required the one-sided upper bound to show a shortfall
strictly greater than 5\%, and that bound corresponded to 4.81\%. A separate
tokenizer-sensitivity criterion passed because Polyglot-Ko assigned the same
Korean strings 13.3\% fewer token units. Tokenizer choice therefore
demonstrably changes token counts, but this study did not establish the
prespecified efficiency magnitude or identify what causal share of the
Korean token surplus belongs to translation, model fit, or tokenizer design.

\Needspace{8\baselineskip}
\section*{Reproducibility}

Code, tests, the frozen specification, and confirmatory analysis outputs
(bootstrap replicates, provenance, and per-pair measurements with corpus text
removed) are in the
\href{https://github.com/damifiance/korean-token-tax/tree/2ebe80ea7dae8722d69cb918d3bbbacbde47e317}{archived
project revision}. The corpus is rebuilt from a hash-verified OPUS archive by
a script that reproduces the split. The estimator snapshot and baseline
tokenizer asset are verified before scoring; all resolved tokenizer files and
revisions are recorded in the run provenance. The historical-tokenizer
comparison was computed post hoc and is not part of the archived
confirmatory outputs.

\bibliographystyle{plainnat}
\bibliography{refs}

\appendix

\section{Pilot precision simulation}

Before the thresholds were chosen, the pilot's 69 per-article sums were
resampled with replacement to 278 articles 400 times; within each simulated
sample a 2{,}000-replicate cluster bootstrap produced the bounds that the
confirmatory rules use. Table~\ref{tab:precision} reports the quantiles of
those bounds across simulations. The observed upper bound, 0.9519, lies near
the upper tail (about the 90th empirical percentile) of this simulated
distribution and below its 95th percentile, 0.9530.

\begin{table}[H]
\centering\small
\caption{Quantiles across 400 simulated held-out samples (278 articles) of
the decision bounds, using pilot per-article sums as the population.}
\label{tab:precision}
\begin{tabular}{@{}lccccc@{}}
\toprule
Bound & 5\% & 25\% & 50\% & 75\% & 95\% \\
\midrule
$\RT$ lower 95\%    & 1.3325 & 1.3443 & 1.3518 & 1.3586 & 1.3684 \\
$\RB$ 90\% CI low   & 1.2563 & 1.2654 & 1.2723 & 1.2782 & 1.2865 \\
$\RB$ 90\% CI high  & 1.2863 & 1.2947 & 1.3013 & 1.3067 & 1.3157 \\
$\Reta$ upper 95\%  & 0.9399 & 0.9440 & 0.9467 & 0.9495 & 0.9530 \\
$\CK$ upper 95\%    & 0.8853 & 0.8902 & 0.8933 & 0.8970 & 0.9030 \\
\bottomrule
\end{tabular}
\end{table}

\section{Exclusion sensitivity on the pilot}

Before the exclusion rules were chosen, every candidate rule was applied to
the pilot. Dropping identical-text rows, rows without Hangul, rows shorter
than five characters, rows with a Korean/English character ratio outside
$[0.2,1.5]$, or all of these together moved $\Reta$ within 0.938--0.941 and
$\RT$ within 1.368--1.393. The pooled estimands are insensitive to these
choices; the rules adopted remove untranslated material on grounds of
validity rather than effect.

\end{document}
